\section{A constructive gap around the natural-number lattice}
\label{sec:gap}

We now give the analytic part of the reduction. Its classical ingredients are
Richardson's lattice penalty and a dense elementary curve; see
\cite{richardson,poonen}. Every estimate and numerical bound needed below is
derived explicitly. In particular, our rounding rule acts on rational
approximations, and our density argument produces a rational witness by bounded
search.

\subsection{The polynomial and its majorant}

Let
\[
 P(X_1,\ldots,X_n)=\sum_\alpha c_\alpha X^\alpha
 \quad(c_\alpha\in\Z),\qquad n\ge1.
\]
The decision problem concerns roots in $\N^n$. A dummy variable may be added
when necessary, so the condition $n\ge1$ loses nothing. Define the positive
integers
\[
 C=\max\left(1,\sum_\alpha\abs{c_\alpha}\right),\qquad
 D=\max\left(1,\deg P\right),
\]
with degree zero assigned to the zero polynomial for this convention. For a
tuple $u=(u_1,\ldots,u_n)$ put
\begin{align}
 R(u)&=1+\sum_{i=1}^n u_i^2,
 &L(u)&=CD\,[2R(u)]^D,\label{eq:majorant}\\
 H_P(u)&=P(u)^2+L(u)^2\sum_{i=1}^n\sin^2(\pi u_i),
 &F_P(u)&=4H_P(u)-1.\label{eq:gap-function}
\end{align}
These are finite elementary expressions. The number $D$ is a fixed exponent
once the input polynomial is given; its power is expanded by repeated
multiplication in the expression syntax. We will apply $F_P$ only to
nonnegative tuples. Later, squaring the coordinates of a dense curve will
enforce this condition without importing a theorem about sums of squares.

\subsection{Rounding without a real comparison}

\begin{lemma}[Coarse rational rounding]\label{lem:rounding}
Suppose $r\in\Q$ approximates a real quantity $x$ with
$\abs{x-r}\le1/8$. Set
\[
 z=\lfloor r+1/2\rfloor\in\Z.
\]
Then
\[
 \abs{x-z}\le5/8,
 \qquad
 \abs{x-z}\le\abs{\sin(\pi x)}.
\]
If $x\ge0$, the computed integer $z$ is nonnegative.
\end{lemma}

\begin{proof}
Rational floor is an exact integer-division operation. It gives
$\abs{r-z}\le1/2$, and consequently
\[
 \abs{x-z}\le\abs{x-r}+\abs{r-z}\le5/8.
\]
For $0\le a\le5/8$, the elementary enclosures
$3<\pi<16/5$ imply $0\le\pi a\le2$. The alternating sine series gives
$\sin s\ge s-s^3/6$ on $[0,2]$. Therefore
\[
 \sin(\pi a)
 \ge\pi a\left(1-\frac{(\pi a)^2}{6}\right)
 \ge\frac{\pi a}{3}\ge a.
\]
Apply this to $a=\abs{x-z}$. Sine's parity and its change of sign on
translation by an integer multiple of $\pi$ show that
$\abs{\sin(\pi x)}=\sin(\pi a)\ge a$. Finally, if $x\ge0$, then
$r\ge-1/8$, whence $r+1/2\ge3/8$ and its rational floor is nonnegative.
\end{proof}

This lemma does not ask whether a real number lies above or below a
half-integer. Even at a half-integer the rational rule is well defined, and
its stated error estimate suffices. The numerical bounds on $\pi$, the sine
series estimate, and the addition identities have finite rational
approximation proofs; explicit approximation primitives are given in the
appendix.

\subsection{A direct perturbation estimate}

\begin{lemma}[Polynomial perturbation]\label{lem:perturbation}
If $\abs{u_i-z_i}\le1$ for every $i$ and
$\delta=\max_i\abs{u_i-z_i}$, then
\[
 \abs{P(u)-P(z)}\le L(u)\delta.
\]
\end{lemma}

\begin{proof}
Since $R(u)\ge1+u_i^2\ge\abs{u_i}$ and $R(u)\ge1$,
\[
 \abs{z_i}\le\abs{u_i}+1\le2R(u).
\]
For a monomial of total degree $d$, regard its value as a product of $d$
factors, allowing repetitions of a coordinate. The telescoping identity
\[
 \prod_{j=1}^d a_j-\prod_{j=1}^d b_j
 =\sum_{j=1}^d
   \left(\prod_{\ell<j}b_\ell\right)
   (a_j-b_j)
   \left(\prod_{\ell>j}a_\ell\right)
\]
bounds its change by $d[2R(u)]^{d-1}\delta$. For $1\le d\le D$ this is at
most $D[2R(u)]^D\delta$. Constant monomials have zero change. Multiplying
by the absolute values of the coefficients and summing proves the estimate.
\end{proof}

A finite maximum of real quantities is constructively definable using
$\max(a,b)=(a+b+\abs{a-b})/2$. Thus the maximum in this proof is not a
choice principle. Alternatively one can carry a common rational error bound
through the same telescoping calculation. Neither presentation adds a
maximum operation to the compiled expression.

\begin{theorem}[The lattice gap]\label{thm:lattice-gap}
For every nonnegative tuple $u$, the rational rounding construction produces
$z\in\N^n$ satisfying
\begin{equation}
 P(z)^2\le2H_P(u).\label{eq:integer-gap-bound}
\end{equation}
Consequently:
\begin{enumerate}
\item if $P$ has no root in $\N^n$, then $F_P(u)\ge1$ for every
      nonnegative $u$;
\item if $F_P(u)<1$, the constructed tuple $z$ is a natural-number root
      of $P$;
\item at every natural-number root $m$, one has $F_P(m)=-1$.
\end{enumerate}
\end{theorem}

\begin{proof}
Approximate each $u_i$ by a rational to error at most $1/8$ and apply
Lemma~\ref{lem:rounding}. The integers obtained are nonnegative and satisfy
\[
 \delta^2:=\max_i\abs{u_i-z_i}^2
 \le\sum_i\sin^2(\pi u_i).
\]
By Lemma~\ref{lem:perturbation} and $(a+b)^2\le2a^2+2b^2$,
\[
 P(z)^2
 \le2P(u)^2+2\abs{P(z)-P(u)}^2
 \le2P(u)^2+2L(u)^2\delta^2
 \le2H_P(u).
\]
If there is no root, the particular integer $P(z)$ is nonzero. Integer
comparison is decidable, and a nonzero integer has square at least one.
Thus $H_P(u)\ge1/2$, giving $F_P(u)\ge1$. Conversely, $F_P(u)<1$ means
$H_P(u)<1/2$. Equation~\eqref{eq:integer-gap-bound} then puts the
nonnegative integer $P(z)^2$ below one, forcing $P(z)=0$. At a root $m$,
both $P(m)$ and every $\sin(\pi m_i)$ vanish, so $H_P(m)=0$ and
$F_P(m)=-1$.
\end{proof}

The first conclusion is a uniform gap, not just pointwise positivity. In
the second conclusion the root is computed from coarse approximations to
the given tuple. No search through all Diophantine tuples is involved.

The growing coefficient $L(u)^2$ has an essential role. For example,
$P(X,Y)=2XY-1$ has no natural-number root, since $2XY$ is even. Yet at
$u=N$, $v=1/(2N)$, with $N\ge1$ an integer, one has $P(u,v)=0$ and
\[
 \sin^2(\pi u)+\sin^2(\pi v)
 =\sin^2\!\left(\frac{\pi}{2N}\right),
\]
which becomes arbitrarily small. Thus the unweighted expression
$4(P(u,v)^2+\sin^2(\pi u)+\sin^2(\pi v))-1$ becomes negative despite
the absence of a root. Near the lattice, the polynomial can vary rapidly
when another coordinate is large. The explicit majorant compensates for
exactly this effect: it bounds the change in $P$ between the tested point
and the rounded lattice point. Merely adding a fixed sine penalty would
not give the reduction.

\subsection{A specified neighborhood of each root}

We also need a rational radius on which the negative value survives. Given
a natural-number root $m$, define
\begin{equation}
 \begin{gathered}
 M=\max_i m_i,\qquad B_0=M+1,\\
 L_0=CD\,[2(1+nB_0^2)]^D,
 \qquad \rho=\frac{1}{16(n+1)L_0}.
 \end{gathered}\label{eq:root-radius}
\end{equation}

\begin{lemma}[Explicit negative neighborhood]\label{lem:root-neighborhood}
If $u_i\ge0$ and $\abs{u_i-m_i}\le\rho$ for every $i$, then
$F_P(u)\le-3/4$.
\end{lemma}

\begin{proof}
We have $\rho\le1$, so $0\le u_i\le B_0$ and $L(u)\le L_0$.
The same monomial telescoping calculation, now on the box $[0,B_0]^n$,
gives $\abs{P(u)}=\abs{P(u)-P(m)}\le L_0\rho$.
By sine's Lipschitz estimate and integer periodicity,
\[
 \abs{\sin(\pi u_i)}
 =\abs{\sin(\pi(u_i-m_i))}\le\pi\rho\le4\rho.
\]
It follows that
\[
 H_P(u)\le L_0^2(1+16n)\rho^2
 =\frac{1+16n}{256(n+1)^2}\le\frac1{16},
\]
where the last inequality uses $1+16n\le16(n+1)^2$. Hence
$F_P(u)=4H_P(u)-1\le-3/4$.
\end{proof}

The proof supplies a radius by integer and rational arithmetic. An
unquantified appeal to continuity would not supply the witness transformation
needed later; equation~\eqref{eq:root-radius} does.
